Un bloc de massa 2 kg, que inicialment es troba en repòs, està lligat a l'extrem d'una
molla de constant elàstica 300 N/m i de longitud natural 30 cm. Per l'altre extrem, la
molla està unida a una paret. Desplacem el bloc cap a la dreta fins que la molla assoleix
una longitud total de 45 cm i el deixem anar, de manera que el bloc es posa a oscil·lar al
voltant de la posició d'equilibri. La fricció entre el terra i el bloc és negligible.

\figura[0.5]{fig1.png}

\begin{apartat}{1,25}
Determineu l'amplitud i el període de l'oscil·lació. Escriviu l'equació del moviment.
\end{apartat}

\begin{apartat}{1,25}
Representeu en la quadrícula adjunta l'evolució de l'energia mecànica, de l'energia
cinètica i de l'energia potencial en funció del temps durant un període.
\end{apartat}
